\begin{ex}
Tính: $354 + 215 = \dots$
\shortans{569}
\loigiai{
$354 + 215 = 569$.
}
\end{ex}

\begin{ex}
Tính: $687 - 342 = \dots$
\shortans{345}
\loigiai{
$687 - 342 = 345$.
}
\end{ex}

\begin{ex}
Tính: $3 \times 7 + 120 = \dots$
\shortans{141}
\loigiai{
$21 + 120 = 141$.
}
\end{ex}

\begin{ex}
Tính: $4 \times 8 - 15 = \dots$
\shortans{17}
\loigiai{
$32 - 15 = 17$.
}
\end{ex}

\begin{ex}
Tính: $36 : 4 + 45 = \dots$
\shortans{54}
\loigiai{
$9 + 45 = 54$.
}
\end{ex}

\begin{ex}
Tính: $27 : 3 + 128 = \dots$
\shortans{137}
\loigiai{
$9 + 128 = 137$.
}
\end{ex}

\begin{ex}
Tìm số lớn nhất trong các số: $879, 897, 789, 798$.
\shortans{897}
\loigiai{
897 là số lớn nhất.
}
\end{ex}

\begin{ex}
Tìm số bé nhất trong các số: $512, 521, 502, 520$.
\shortans{502}
\loigiai{
502 là số bé nhất.
}
\end{ex}

\begin{ex}
Số liền trước của 1000 là số nào?
\shortans{999}
\loigiai{
$1000 - 1 = 999$.
}
\end{ex}

\begin{ex}
Số liền sau của 689 là số nào?
\shortans{690}
\loigiai{
$689 + 1 = 690$.
}
\end{ex}

\begin{ex}
Tìm $x$, biết: $x + 234 = 568$.
\shortans{334}
\loigiai{
$x = 568 - 234 = 334$.
}
\end{ex}

\begin{ex}
Tìm $x$, biết: $x - 145 = 320$.
\shortans{465}
\loigiai{
$x = 320 + 145 = 465$.
}
\end{ex}

\begin{ex}
Tìm $x$, biết: $4 \times x = 36$.
\shortans{9}
\loigiai{
$x = 36 : 4 = 9$.
}
\end{ex}

\begin{ex}
Tìm $x$, biết: $x : 3 = 8$.
\shortans{24}
\loigiai{
$x = 8 \times 3 = 24$.
}
\end{ex}

\begin{ex}
Một trường tiểu học có 420 học sinh nữ và 390 học sinh nam. Hỏi trường đó có tất cả bao nhiêu học sinh?
\shortans{810}
\loigiai{
$420 + 390 = 810$ học sinh.
}
\end{ex}

\begin{ex}
Một cửa hàng có 500 kg đường, đã bán 235 kg. Hỏi cửa hàng còn lại bao nhiêu ki-lô-gam đường?
\shortans{265}
\loigiai{
$500 - 235 = 265$ kg.
}
\end{ex}

\begin{ex}
Mỗi bàn học có 2 bạn ngồi. Hỏi 8 bàn học như thế có tất cả bao nhiêu bạn?
\shortans{16}
\loigiai{
$2 \times 8 = 16$ bạn.
}
\end{ex}

\begin{ex}
Có 28 quả cam chia đều vào 4 giỏ. Hỏi mỗi giỏ có bao nhiêu quả cam?
\shortans{7}
\loigiai{
$28 : 4 = 7$ quả cam.
}
\end{ex}

\begin{ex}
Một tam giác có độ dài 3 cạnh đều bằng 8 cm. Tính chu vi tam giác đó theo đơn vị xăng-ti-mét.
\shortans{24}
\loigiai{
$8 \times 3 = 24$ cm.
}
\end{ex}

\begin{ex}
Một hình tứ giác có độ dài 4 cạnh đều bằng 5 cm. Tính chu vi tứ giác đó theo đơn vị xăng-ti-mét.
\shortans{20}
\loigiai{
$5 \times 4 = 20$ cm.
}
\end{ex}

\begin{ex}
Tính: $2 \times 9 + 4 \times 5 = \dots$
\shortans{38}
\loigiai{
$18 + 20 = 38$.
}
\end{ex}

\begin{ex}
Tính: $30 : 5 + 24 : 3 = \dots$
\shortans{14}
\loigiai{
$6 + 8 = 14$.
}
\end{ex}

\begin{ex}
Số tròn chục lớn nhất có ba chữ số là số nào?
\shortans{990}
\loigiai{
Số đó là 990.
}
\end{ex}

\begin{ex}
Số gồm 5 trăm và 7 đơn vị viết là số nào?
\shortans{507}
\loigiai{
Số đó là 507.
}
\end{ex}

\begin{ex}
Một sợi dây dài 1 m, người ta cắt bớt đi 30 cm. Đoạn dây còn lại dài bao nhiêu xăng-ti-mét?
\shortans{70}
\loigiai{
$1\text{ m} = 100\text{ cm}$; $100 - 30 = 70$ cm.
}
\end{ex}
